\section{Levi-Civita connection}

Assume $M$ is a Riemannian manifold.

\begin{definition}[Levi-Civita connection]\label{def:5.1}
\[
\nabla : \underset{\in (X,Y)}{\mathcal{\mathfrak{X}(M)}\times\mathfrak{X}(M)} \to \underset{\in \nabla \times Y}{\mathfrak(M)}
\]
is called the \emph{covariant derivative} or \emph{Levi-Civita connection} if
$\forall X, Y, Z \in \mathfrak{X}(M)$ and $\forall f \in C^{\infty}(M)$:
\begin{enumerate}[label=(\arabic*), ref=Definition \thedefinition~(\arabic*)]
    \item \label{def:5.1.1} $\nabla_{X+Y}Z = \nabla_X Z + \nabla_Y Z$
    \item \label{def:5.1.2} $\nabla_{fX}Y = f\nabla_X Y$
    \item \label{def:5.1.3} $\nabla_X(Y+Z) = \nabla_X Y + \nabla_X Z$
    \item \label{def:5.1.4} $\nabla_X(fY) = X(f)Y + f\nabla_X Y$
    \item \label{def:5.1.5} $\nabla_X Y - \nabla_Y X = [X,Y]$ (torsion-free)
    \item \label{def:5.1.6} $X\langle Y, Z \rangle = \langle \nabla_X Y, Z \rangle + \langle Y, \nabla_X Z \rangle$ (metric compatibility)
\end{enumerate}
\end{definition}

\begin{theorem}[Fundamental theorem of Riemannian geometry]
There exists a unique ($\exists !$) Levici-Civita connection $\nabla$ on a Riemannian manifold $M$.
\end{theorem}

\begin{proof}\label{proof:fundamental_theorem_of_riemannian_geometry}
We start by showing uniqueness. By \autoref{def:5.1.6}:
\[
\begin{aligned}
    X\langle Y, Z \rangle &= \langle \nabla_X Y, Z \rangle + \langle Y, \nabla_X Z \rangle\\
    Y \langle Z, X \rangle &= \langle \nabla_Y Z, X \rangle + \langle Z, \nabla_Y X \rangle\\
    - Z \langle X, Y \rangle &= - \langle \nabla_Z X, Y \rangle - \langle X, \nabla_Z Y \rangle
\end{aligned}
\]
Summing over these and applying \autoref{def:5.1.5}:
\begin{align}\label{eq:levi-civita_uniqueness}
\langle \nabla_XY, Z\rangle &= \frac12 \{X\langle Y, Z \rangle + Y \langle Z, X \rangle - Z \langle X, Y \rangle + \langle [X,Y], Z \rangle \notag \\
& - \langle [Y,Z], X \rangle + \langle [Z,X], Y \rangle\}
\end{align}
This shows uniqueness.

We now show existence. Let the right-hand side of \eqref{eq:levi-civita_uniqueness} be $F_{X,Y}(Z)$.
Then $\forall f_1, f_2 \in C^{\infty}(M)$ and $\forall Z_1, Z_2 \in \mathfrak{X}(M)$:
\[
F_{X, Y}(f_1 Z_1 + f_2Z_2) = f_1F_{X,Y}(Z_1) + f_2F_{X,Y}(Z_2),
\]
is easily verified. For a chart $(U, \phi)$, we can define
\[
\nabla_X^{(U,\phi)}Y := \sum_{i, j = 1}^{n}F_{X,Y}(\partial_i)g^{ij}\partial_j,
\]
where $(g^{ij})=(g_{ij})^{-1}$. For $Z = \sum_{k=1}^n Z^k\partial_k$, using that $\sum_{j=1}^n g^{ij}g_{jk}=\delta^i_k$, we have
\[
\begin{aligned}
\langle \nabla_X^{(U,\phi)}Y, Z \rangle &= \sum_{i,j,k = 1}^n F_{X,Y}(\partial_i)g^{ij}Z^k g_{jk} \\
&= \sum_{i=1}^n F_{X,Y}(\partial_i)Z^i = F_{X,Y}(Z).
\end{aligned}
\]
Note that this is independent of the choice of chart $(U,\phi)$! This implies that for any other chart $(V, \psi)$ such that $U \cap V \neq \emptyset$, $\nabla_X^{U,\phi} = \nabla_X^{(V,\psi)}Y$.
Thus $\nabla_X Y := \nabla_X^{(U, \phi)}Y$ yields a globally defined vector field. Direct calculation verifies all axioms \autoref{def:5.1.1}--\autoref{def:5.1.6} of \autoref{def:5.1}.
\end{proof}

\begin{problem}
Verify axioms \autoref{def:5.1.1}--\autoref{def:5.1.6} of \autoref{def:5.1} for the Levi-Civita connection $\nabla$ constructed in \autoref{proof:fundamental_theorem_of_riemannian_geometry}.
\end{problem}

\begin{definition}[Christoffel symbols]\label{def:christoffel_symbols}
For a chart $(U,\phi)$, we define functions $\Gamma_{ij}^k$ by
\[
\nabla_{\partial_i}\partial_j = \sum_{k=1}^n \Gamma_{ij}^k \partial_k
\]
called \emph{Christoffel symbols}.
\end{definition}

\begin{remark}
$\nabla$ is NOT a tensor field, implying that $\Gamma_{ij}^k$ are not the components of a tensor.
\end{remark}

\begin{proposition}\label{prop:christoffel_symbols}
$\forall i, j, k = 1, \dots, n$:
\begin{enumerate}[label=(\arabic*), ref=Proposition \theproposition~(\arabic*)]
    \item \label{item:christoffel_1} $\Gamma_{ij}^k = \sum_{l=1}^n \langle \nabla_{\partial_i}\partial_j, \partial_l\rangle g^{lk} = \frac12 \sum_{l=1}^n g^{kl}(\partial_ig_{jl}+\partial_jg_{il}- \partial_lg_{ij})$.
    \item \label{item:christoffel_2} $\Gamma_{ij}^k = \Gamma_{ji}^k$.
    \item \label{item:christoffel_3} $\partial_kg_{ij} = \frac{\partial g_{ij}}{\partial x^k} = \sum_{l=1}^n (\Gamma_{ik}^lg_{lj} + \Gamma_{jk}^lg_{li})$.
\end{enumerate}
\end{proposition}

\begin{proof}
By \autoref{def:christoffel_symbols}: $\langle \nabla_{\partial_i}\partial_j, \partial_l \rangle = \Gamma_{ij}^k g_{kl}$.
Multiplying by $g^{lm}$ and summing over $l$ yields the first equality \autoref{item:christoffel_1}.
The second equality follows from \autoref{def:5.1} with
$X = \partial_i$, $Y = \partial_j$, $Z = \partial_l$ and $[\partial_i, \partial_j]$.

\autoref{item:christoffel_2} follows from \autoref{item:christoffel_1} and the symmetry of $g_{ij}$.

\autoref{item:christoffel_3} follows from \autoref{def:5.1} with $X = \partial_k$, $Y = \partial_i$, $Z = \partial_j$ and \autoref{item:christoffel_2}.
\end{proof}

For local expressions $X = \sum_{i=1}^nX^i\partial_i$ and $Y = \sum_{j=1}Y^j\partial_j$, calculation using \autoref{def:5.1} \autoref{def:5.1.1}--\autoref{def:5.1.4} yields
\[
\begin{aligned}
\nabla_XY &= \sum_{i=1}^nX^i\nabla_{\partial_i}\left(\sum_{j=1}^nY^j\partial_j\right) \\
&= \sum_{i=1}^nX^i\left(\sum_{j=1}^n(\partial_iY^j \partial_j + Y^j\sum_{k=1}^n\Gamma_{ij}^k\partial_k)\right) \\
&= \sum_{k=1}^n \left\{X(Y^k) + \sum_{i,j=1}^n \Gamma_{ij}^kX^iY^j\right\}\partial_k.
\end{aligned}
\]
A key fact is that $(\nabla_XY)_p$ depends only on $X_p$ and the behaviour of $Y$ along $X_p$ direction at $p$.
Thus, for $v \in T_pM$, we define
\[
\nabla_vY := (\nabla_XY)_p \text{ such that } X_p = v,
\]
which is well-defined and independent of the extension $X$.

\begin{lemma}\label{lemma:christoffel_symbol_transformation}
Let $(U,\phi)$ and $(V, \psi)$ be charts such that $U \cap V \neq \emptyset$.
Let $\phi = (x^1, \dots, x^n)$ and $\psi = (y^1, \dots, y^n)$ and $\Gamma_{ij}^k, \hat{\Gamma}_{ij}^k$
be Chistoffel symbols with respect to each chart. Then $\forall \alpha, \beta, \gamma = 1, \dots, n$:
\[
\hat{\Gamma}_{\alpha\beta}^{\gamma} = \sum_{k=1}^n \left( \frac{\partial^2x^k}{\partial y^{\alpha}\partial y^{\beta}} + \sum_{i,j=1}^n \frac{\partial x^i}{\partial y^{\alpha}}\frac{\partial x^j}{\partial y^{\beta}}\Gamma_{ij}^k \right)\frac{\partial y^{\gamma}}{\partial x^k}.
\]
\end{lemma}

\begin{proof}
By the transformation law $\frac{\partial}{\partial y^{\alpha}} = \sum_{i=1}^n \frac{\partial x^i}{\partial y^{\alpha}}\frac{\partial}{\partial x^i}$:
\[
\nabla_{\frac{\partial}{\partial y^{\alpha}}} \frac{\partial}{\partial y^{\beta}} = \sum_{i,j=1}^n \left( \frac{\partial x^i}{\partial y^{\alpha}} \left(\frac{\partial}{\partial x^i}\frac{\partial x^j}{\partial y^{\beta}}\right) \frac{\partial}{\partial x^j} + \frac{\partial x^i}{\partial y^{\alpha}}\frac{\partial x^j}{\partial y^{\beta}}\sum_{k=1}^n\Gamma_{ij}^k\frac{\partial}{\partial x^k}\right).
\]
Using $\frac{\partial}{\partial x^i} = \sum_{\gamma = 1}^n \frac{\partial y^{\gamma}}{\partial x^i}\frac{\partial}{\partial y^{\gamma}}$:
\[
\sum_{i=1}^n \frac{\partial x^i}{\partial y^{\alpha}}\left(\frac{\partial}{\partial x^i}\frac{\partial x^j}{\partial y^{\beta}}\right) = \sum_{i,\gamma=1}^n \frac{\partial x^i}{\partial y^{\alpha}}\frac{\partial y^{\gamma}}{\partial x^i}\frac{\partial^2x^j}{\partial y^{\gamma}\partial y^{\beta}} = \frac{\partial^2x^j}{\partial y^{\alpha}\partial y^{\beta}}.
\]
Note that $\frac{\partial}{\partial x^i}\frac{\partial x^j}{\partial y^{\beta}} = \frac{\partial}{\partial y^{\beta}}\left(\frac{\partial x^j}{\partial x^i}\right) = \frac{\partial}{\partial y^{\beta}}\delta_i^j = 0$.
Thus,
\[
\begin{aligned}
\sum_{\gamma = 1}^n \hat{\Gamma}_{\alpha\beta}^{\gamma}\frac{\partial}{\partial y^{\gamma}}
&= \nabla_{\frac{\partial}{\partial y^{\alpha}}} \frac{\partial}{\partial y^{\beta}} \\
&= \sum_{k=1}^n \left(\frac{\partial^2x^k}{\partial y^{\alpha}\partial y^{\beta}} + \sum_{i,j=1}^n \frac{\partial x^i}{\partial y^{\alpha}}\frac{\partial x^j}{\partial y^{\beta}}\Gamma_{ij}^k \right)\frac{\partial}{\partial x^k} \\
&= \sum_{\gamma=1}^n\sum_{k=1}^n \left(\frac{\partial^2x^k}{\partial y^{\alpha}\partial y^{\beta}} + \sum_{i,j=1}^n \frac{\partial x^i}{\partial y^{\alpha}}\frac{\partial x^j}{\partial y^{\beta}}\Gamma_{ij}^k \right)\frac{\partial y^{\gamma}}{\partial x^k}\frac{\partial}{\partial y^{\gamma}}.
\end{aligned}
\]
\end{proof}